\begin{chapterbox}
    \chapter{In vertrauten Landen (WIP)}
    \label{wip-in-vertrauten-landen}
    \az{Jahr 75}
    
    \begin{center}
        Ein Neues Abenteuer von Bragor, dem Tarus aus dem Sturmtal

        Teil C der Neuen Abenteuer

        Fortsetzung des Zwischenspiels \hyperref[wip-zwei-folgenschwere-tode]{Zwei folgenschwere Tode (WIP)} und von \hyperref[wip-kampf-um-die-flusslande]{Kampf um die Flusslande (WIP)}
    \end{center}

    Sina und Patia fliehen von Andor nach Sturmtal. Dort leidet Bragor unter Verlusten.
\end{chapterbox}




Diese Fan-Geschichte ist noch in Arbeit. Ich teile sie bereits hier, damit sie für mich nicht verloren geht, falls mein Laptop spontan verschwinden sollte, und damit sie für euch nicht völlig verloren geht, falls ich spontan verschwinden sollte. Fühlt euch in letzterem Fall vollkommen frei, hierrin zu schmökern und nach Lust und Laune weiter daran zu werkeln. Aber solange ich ein aktiver Andori bleibe, lest bitte nicht zu genau nach, was ich geplant haben könnte. LG BBB :D



\newpage
\section{Prolog}
\az{Jahr 74}

Sina und Patia erlebten eine unangenehme Begegnung. Patia beschloss, Andor zu verlassen und ins Sturmtal zu ziehen.





\newpage
\section{Kapitel}
\az{Jahr 75}

Sina und Patia hatten sich ins Sturmtal eingelebt. 

Bragor litt weiterhin stark unter den Verlusten, die die Helden in Hadria und beim Untergang der Aldebaran erlitten hatten.

Die Schwarze Kogge griff Sturmtal an und stahl uralte \gend{Druid*innen}{Druiden}artefakte.

Sina wurde entführt.

Bragor und Rhega begriffen, dass die Anwesenheit der Schwarzen Kogge hier bedeutete, dass Varatans Fluch gebrochen worden war. Vielleicht waren die auf Narkon gekenterten \gend{Held*innen}{Helden} nicht verloren. Doch konnte Bragor sich nicht dazu aufraffen, Sturmtal zu verlassen.

Rhega und Patia reisten nach Werftheim und vernahmen vor der Taverne zur Tanzenden Tare eine Gruppe \gend{Bard*innen}{Barden}, die die Geschichte vom Bruch des Fluchs zum Besten gaben.



\newpage
\section{Kapitel}

Rhega reiste alleine nach Sturmtal zurück, um Bragor die guten Neuigkeiten zu berichten. 

Gemeinsam fanden die beiden heraus, dass die Kogge-Mannschaft ein Ritual an Varatans Leuchtturm durchzuführen plante.

Rhega und Bragor reisten per Boot zur Ruinenstadt. 

Sie fanden die Besatzung der Kogge, doch wurden rasch selbst entdeckt.

Bragor und Rhega wurden von Kentars rotem Stein unglaublich schnell und stark überwältigt, als hätten die beiden eine besondere Schwäche dafür. 

Rhega versuchte sich mit \gend{Druid*innen}{Druiden}kraft zu wehren, wurde als Thoggers Enkelin und mächtige Druidin erkannt und entführt.

Bragor konnte knapp fliehen, weil keine Aufmerksamkeit auf ihm Nicht-Druiden lag.

Er sah mit an, wie Ean Quella in der Ferne etwas Magisches tat, um Varatans Leuchtturm erscheinen zu lassen.



\newpage
\section{Kapitel}

Patia reiste alleine nach Silberhall weiter. 

Auf der Bootsstrecke traf sie Barz wieder, der ebenfalls die Hoffnung hatte, in Silberhall alle \gend{Held*innen}{Helden} vorzufinden. 

Diese Hoffnung wurde etwas gedämpft, als die beiden in Silberhall nur die frisch aus Narkon befreiten Iril und Tenaya wiederfanden. 

Die restlichen \gend{Held*innen}{Helden} waren offenbar wieder im nördlichen Teil Silberlands unterwegs, um Überlebende des Fluches zu versammeln und in Sicherheit zu führen.


\newpage
\section{Kapitel}

Bragor verzweifelte in der Ruinenstadt an den Überbleibseln von Kentars Manipulationen, bis er einen emotionalen Durchbruch hatte.

Er fand im tiefen Schnee, der mit der Kogge gekommen war, ein magisches Portal, wie es sie vor allem in Hadria gab.

Bragor konnte durch das Portal einen Hilferuf nach Silberhall schicken.

Die dort anwesenden \gend{Held*innen}{Helden} erreichten ihn gerade noch, bevor das Portal sich schloss.

Tenaya, die schon einmal an Varatans Leuchtturm vorbeigesegelt war\sidenote{\siehechronik{Tenaya, die Wächterin des Feuers (2022)}}, fand heraus, wie man ihn erscheinen ließ.

Das Quartett machte sich bereit, sich der Schwarzen Kogge zu stellen.

\newpage
\section{Kapitel}

Rhega wurde gezwungen, mit Thogger und den gestohlenen Artefakten ein Ritual an Varatans Leuchtturm zu starten.

Mittendrin konnte Rhega mit Sina fliehen.

In den Klippen um Varatans Leuchtturm fand Rhega die restlichen \gend{Held*innen}{Helden}, die mitten in einer Besprechung waren, wie man die beiden am besten befreite.

Die \gend{Held*innen}{Helden} versuchten, die Schwarzen Kogge fahrtüchtig zu machen und zu fliehen, während die Besatzung auf dem Leuchtturm war. 

Doch die Besatzung folgte ihnen an Bord. Es kam zum Gefecht.

Erneut wurden Bragor, Rhega und Sina in einer Vision von Kentars rotem Stein gefangen.

Doch dieses Mal half ihnen Thogger, daraus zu entkommen. Dies war sein Loch, nicht ihres.

Nach und nach wurden die \gend{Held*innen}{Helden} über Bord gedrängt. Rhega riss Ean Quella mit sich, ohne deren Magie die Besatzung nicht mehr auf Varatans Turm kommen konnte.

Als nur noch Bragor übrig war, stellte Callem sich ihm und stach ihm sein Schwert durch die Brust.

Bragors Gesichtfeld verdunkelte sich.

Ein ferner Gong ertönte.








\newpage
\section{Epiloge}

Rhega, Tenaya, Iril und Barz kehrten geschlagen nach Silberhall zurück. 

Sina und Patia wurden wieder vereint, Ean Quella in die Aufsicht der \gend{Silberzwerg*innen}{Silberzwerge} zur Heilung gegeben.

In der Silbermine fanden sie die zurückgekehrten Chada, Thorn, Stinner und Kirr vor, ausgerüstet mit zwei Magischen Waffen, dem Sturmschild und einem hadrischen Stundenglas. 

Sie beschlossen gemeinsam, Bragors Tod zu rächen und Jagd auf die Kogge zu machen.

\begin{commentbox}
    Hier wäre eine gute Gelegenheit, wieder einmal die Bonus-Legende \siehechronik{Die Rückkehr der Schwarzen Kogge (2015)} zu spielen.
\end{commentbox}

\trennzeichen

Der frisch aus Narkon befreite Ijsdur erreichte den Nestbaum der Takuri, um Aćh von der Rückkehr nach Andor zu überzeugen.

\trennzeichen

Fenn saß im Kerker unter der Rietburg.

\begin{commentbox}
    Weiter geht es in \hyperref[wip-eine-faehrtenleserin-entscheidet-sich]{Eine Fährtenleserin entscheidet sich (WIP)}.
\end{commentbox}

